\documentclass[11pt,a4paper]{article}
\usepackage[T1]{fontenc}
\usepackage[utf8]{inputenc}
\usepackage{lmodern}
\usepackage[a4paper,margin=25mm]{geometry}
\usepackage{microtype}
\usepackage{amsmath,amssymb}
\usepackage{booktabs,longtable,array}
\usepackage{enumitem}
\usepackage{url}
\usepackage[hidelinks]{hyperref}
\usepackage{tikz}
\usetikzlibrary{arrows.meta,positioning,calc}
\setlength{\parindent}{0pt}
\setlength{\parskip}{6pt}
\setlength{\emergencystretch}{2em}
\renewcommand{\arraystretch}{1.2}
\setlist{nosep,leftmargin=*}
\pagestyle{plain}
\raggedbottom
\date{}
\hypersetup{pdfauthor={Hariprasad},pdftitle={MathStudio - Documentation}}
\begin{document}
\begin{titlepage}\centering
{\large Hariprasad}\\[4pt]
{\large Roll Number: DA24B005}\\[25pt]
{\LARGE\bfseries MathStudio -- Documentation}\\[14pt]
{\large User guide, setup and operations}\\[30pt]
\begin{minipage}{0.85\textwidth}
This document covers using the solver, setting up the project, running the pipeline and monitoring the services. The design, training results and application outputs are in the separate project report.
\end{minipage}
\end{titlepage}
\setcounter{page}{2}
\section*{User Manual}
\subsection*{Open the application}
Once the operator starts the project, open \textbf{http://localhost:8080}. The workspace has a question box, a \textbf{Solve problem} button and three examples. No account is required. The \textbf{Guide} tab provides these instructions inside the application.

\subsection*{Ask a question}
\begin{enumerate}
\item Type one mathematics question. Include the numbers, units and what you want to find.
\item For an explanation, add “show the steps.” For a short response, add “give only the final value.”
\item Select \textbf{Solve problem}, or press \textbf{Ctrl + Enter} (Command + Enter on Mac).
\item Wait for the solution below the question box. The button is disabled while answering. The first answer can take longer while the service starts.
\item Read the answer and check its reasoning. Select \textbf{Copy answer} to save the text elsewhere.
\end{enumerate}
Examples: “Solve 2x + 3 = 11 for x”, “What is 3/4 plus 1/8?”, and “Find the area of a rectangle with length 9 and width 4.” Selecting an example fills the question box; you can edit it before solving.

\subsection*{Recover from a problem}
An empty question produces “Enter a mathematics question first.” Questions are limited to 8,000 characters. If the local service is unavailable, wait a moment and retry; if it remains unavailable, contact the operator. For a timeout, shorten the question. A length-limit notice means the response may be incomplete; ask a more focused question. Equations render locally and work without an external font or script service. The answer is generated by a small model and may be incorrect, especially for difficult or unfamiliar problems.

\subsection*{Inspect the process}
Open \textbf{Pipeline} for the current data count, split sizes, Spark speed, training duration, generation smoke result and model registry decision. Each stage shows its state and duration. The execution log and automatic five-second refresh reveal progress during a run. \textbf{Refresh evidence} requests an immediate update. Tool links open local MLflow, Prometheus and Grafana screens.

A five-question smoke result describes only those five examples. Completion loss measures fit to reference solutions; it is not a mathematical correctness percentage. Missing reports are shown as awaiting evidence. During a new run, existing evaluation files can still show the previous completed evaluation until their stage replaces them.

\subsection*{Privacy}
The deployed inference service runs on this computer. Request logs contain endpoint, outcome and timing, not the question text. The browser holds the current question and answer in the open page. Initial operator setup downloads public data and model artifacts. Do not enter passwords or private credentials into a mathematics question.

Operator installation, restart and troubleshooting instructions are in Operator runbook.

\clearpage
\section*{Monitoring and Incident Operations}
The local stack now has nine services: frontend, backend, MLflow model serving, MLflow tracking, Prometheus, Grafana, Alertmanager, Blackbox Exporter and a persistent alert receiver. Seven exporter targets and eight HTTP endpoint probes are collected every fifteen seconds. Six recording rules and seventeen alert rules are provisioned as source files.

\subsection*{Operator views}
\begingroup\small
\begin{longtable}{@{}>{\raggedright\arraybackslash}p{0.5000\dimexpr\textwidth-12pt\relax}>{\raggedright\arraybackslash}p{0.5000\dimexpr\textwidth-12pt\relax}@{}}
\toprule
\textbf{View} & \textbf{Purpose} \\
\midrule
\endfirsthead
\textbf{View} & \textbf{Purpose} \\
\midrule
\endhead
Grafana Service Operations & Metrics/HTTP readiness, incident count, solve throughput, 5xx ratio, p50/p95/p99 latency, process RSS/CPU, in-flight requests and scrape timing \\
Grafana Alert Center & Severity queue, pending/firing state, notification delivery/failures, heartbeat age, silences and operating policy \\
Grafana ML Pipeline & Last measured stage durations and outcomes, running-stage elapsed time, Spark throughput, five-case acceptance score and model state \\
Alertmanager & Current grouped alerts, inhibition status and bounded maintenance silences \\
Local alert inbox & Delivered firing/resolved event history, heartbeat and controlled-drill history \\
\bottomrule
\end{longtable}
\endgroup
The service dashboard is \url{http://localhost:13000/d/local-math-assistant.} Alert Center is \url{http://localhost:13000/d/math-alert-center;} ML lifecycle is \url{http://localhost:13000/d/math-mlops.} Alertmanager runs at \url{http://localhost:19093/} and the inbox at \url{http://localhost:19199/.}

\subsection*{Alert policy and scope}
Availability alerts hold for one minute. Traffic alerts exclude HTTP 4xx from the server-error numerator. Warning errors exceed 5\% with at least five requests in ten minutes; critical errors exceed 20\% with at least ten requests. Both hold for two minutes. End-to-end solve p95 latency above thirty seconds requires ten requests in ten minutes and a five-minute hold. Additional rules cover in-flight queue pressure, model RSS, pipeline parent/stage failures, stalled training, generation-gate failure, configuration reload failure and notification-delivery failure.

Latency buckets extend to 180 seconds for CPU inference. In-flight counts include waiting requests because inference is serialized. RSS and CPU are process metrics; 100\% CPU means one core. Observed HTTP uptime describes collected samples, not a guaranteed full-day production SLO. No traffic produces no latency estimate. The five-case smoke score remains distinct from broad mathematical correctness.

Prometheus evaluates alerts, and Alertmanager groups, deduplicates, inhibits and routes them. Critical incidents suppress warnings for the same component; scrape outages suppress duplicate endpoint alerts. The local receiver accepts bounded validated notifications, flushes them to a rotating journal, and responds with a retryable error when durable storage is unavailable. The UI renders event text safely. External email/chat notification destinations are not configured.

\newpage
\subsection*{Verification and operations}
The Watchdog alert is intentionally always firing and delivers a heartbeat roughly every minute. It is excluded from incident counts. A missing-heartbeat rule has startup grace. A temporary delivery drill goes through the real metric, rule, routing and webhook path and must deliver both firing and resolved notifications.

\begin{quote}\small
\begin{verbatim}
.venv/bin/python scripts/verify_monitoring.py --drill
\end{verbatim}
\end{quote}
The verifier checks all dashboard PromQL expressions, healthy target inventory, readiness probes, rule health, routing, the persistent inbox and recent Watchdog delivery. Promtool tests cover hold/recovery behavior, low-volume gating, thresholds, lifecycle and 4xx exclusion. Python tests cover durable concurrent storage, bounded input/history, failure retry, drill expiry and the new API/pipeline telemetry. The refreshed suite has 58 passing cases.

Named Docker volumes retain monitoring state across recreation. Prometheus retention is bounded to fifteen days or two GB; inbox files rotate at two MiB with four backups. The UI shows the latest 200 deliveries from the current journal. A timed silence should identify an owner and the smallest affected service. The complete recovery and maintenance guide is \texttt{monitoring/\allowbreak{}RUNBOOK.\allowbreak{}md}, also served from \url{http://localhost:19199/runbook.}

This is a local operations setup. An entire host outage requires monitoring and delivery outside that host. Production exposure would additionally require access control, capacity policy, secret management, backups and deployment-specific threshold calibration.

\clearpage
\section*{Operator Runbook}
For the complete guided setup, use \texttt{.\allowbreak{}/\allowbreak{}scripts/\allowbreak{}bootstrap\_\allowbreak{}project.\allowbreak{}sh}. For an existing environment, use \texttt{.\allowbreak{}/\allowbreak{}scripts/\allowbreak{}ta\_\allowbreak{}evaluate.\allowbreak{}sh} to rerun the pipeline, launch all services and verify the result. Both support \texttt{-\allowbreak{}-\allowbreak{}dry-\allowbreak{}run}; the TA quickstart describes their options and resource links. The individual commands below remain available for diagnosis.

\subsection*{Requirements and local setup}
Use Python 3.14.6, Java 17 or newer for Spark, Docker with Compose, about 8 GiB of available memory for CPU training and several GiB of free disk. The measured run uses four CPU threads. A project-local Java 17 runtime is already available on the demonstration machine; it is intentionally excluded from source control and the submission archive.

For a new machine, unpack the submission archive and run \texttt{.\allowbreak{}/\allowbreak{}scripts/\allowbreak{}setup.\allowbreak{}sh}. Install Java and set \texttt{JAVA\_\allowbreak{}HOME} before running Spark. The pinned dependencies are in \texttt{requirements.\allowbreak{}lock}; \texttt{python\_\allowbreak{}env.\allowbreak{}yaml} and \texttt{MLproject} describe MLprojects execution. The current working environment passes \texttt{pip check}. The first public model/dataset download needs network access; serving is offline once the Qwen cache is populated.

\subsection*{Rebuild the complete selected-model pipeline}
\begin{quote}\small
\begin{verbatim}
.venv/bin/python -m src.pipeline.run --force
\end{verbatim}
\end{quote}
The full Qwen run includes inspection, 800-row ingestion, EDA, Spark cleaning, deterministic splitting, a three-step CPU benchmark, one-epoch LoRA training, base and adapter held-out evaluation, generated-answer acceptance and MLflow registration. The benchmark enforces a six-GiB process RSS limit. On this machine training takes approximately 25-30 minutes. The full data snapshot contains 799 cleaned rows; normalized question deduplication leaves 798 split rows. Fully truncated training examples are counted separately.

For a short recording of data preparation:

\begin{quote}\small
\begin{verbatim}
.venv/bin/python -m src.pipeline.run --target split --force
\end{verbatim}
\end{quote}
This updates data-stage timing reports. For a cached reproducibility check, omit \texttt{-\allowbreak{}-\allowbreak{}force}. Use \texttt{dvc dag} to visualize the actual dependencies. Every run is local. No command pushes source, models or artifacts to GitHub or a cloud platform.

\subsection*{Start and verify the demonstration}
\begin{quote}\small
\begin{verbatim}
./scripts/start_demo.sh
.venv/bin/python scripts/verify_runtime.py
\end{verbatim}
\end{quote}
The start script populates the pinned public model cache if needed and starts the nine Compose services. Local URLs:

\begingroup\small
\begin{longtable}{@{}>{\raggedright\arraybackslash}p{0.5000\dimexpr\textwidth-12pt\relax}>{\raggedright\arraybackslash}p{0.5000\dimexpr\textwidth-12pt\relax}@{}}
\toprule
\textbf{Screen} & \textbf{URL} \\
\midrule
\endfirsthead
\textbf{Screen} & \textbf{URL} \\
\midrule
\endhead
Workspace & \url{http://localhost:8080} \\
Pipeline console & \url{http://localhost:8081/} \\
API specification & \url{http://localhost:8000/docs} \\
MLflow experiments and registry & \url{http://localhost:5000} \\
MLflow model serving & \url{http://localhost:5002/ping} \\
Prometheus & \url{http://localhost:19090} \\
Grafana dashboard & \url{http://localhost:13000/d/local-math-assistant} \\
\bottomrule
\end{longtable}
\endgroup
\texttt{docker compose ps} displays service health. The model service loads the \texttt{champion} registry alias before the backend/frontend readiness dependencies complete. Model exporter metrics are internal to the Compose network. Ports 19090 and 13000 avoid conflicting with another assignment's existing monitoring services.

\subsection*{Restore artifacts and registry in another folder}
\begin{quote}\small
\begin{verbatim}
.venv/bin/dvc pull
./scripts/restore_registry.sh
./scripts/start_demo.sh
\end{verbatim}
\end{quote}
The local DVC remote is \texttt{.\allowbreak{}dvc/\allowbreak{}cache-\allowbreak{}store}, supplied inside the submission archive. MLflow's SQLite history on the demonstration machine contains absolute artifact locations, so do not copy that database into an unrelated path and expect it to work. Restore the adapter and reports with DVC, then run the registration script to create a new local executable model. If new experiment IDs differ, update \texttt{mlflowExperimentId} in \texttt{frontend/\allowbreak{}config.\allowbreak{}js} for direct run links; the original demonstration history is available on its original machine. It verifies base identity, case results and adapter hash. On a restored machine the console labels the submitted stage history as saved evidence; original run IDs are displayed without broken links, while fresh local runs replace that history. A source-only Git clone does not contain ignored artifact objects; use the ZIP, or rebuild the pipeline.

\subsection*{Verification and reports}
\begin{quote}\small
\begin{verbatim}
./scripts/evaluate_submission.sh
./scripts/evaluate_submission.sh --full
./scripts/run_mlproject.sh test
\end{verbatim}
\end{quote}
The fast verification runs tests, shows DVC status/DAG, checks saved acceptance evidence and validates Compose/dashboard configuration. \texttt{-\allowbreak{}-\allowbreak{}full} regenerates all five smoke answers into a separate runtime file. It does not overwrite submitted evaluation evidence. Live HTTP checks are separate and require the stack to be running.

\subsection*{Monitoring and incident response}
Open the Grafana \href{http://localhost:13000/d/math-alert-center}{Alert Center}, the local \href{http://localhost:19199/}{delivery inbox}, or \href{http://localhost:19093/}{Alertmanager} for active alerts and bounded maintenance silences. The detailed operations runbook explains thresholds and recovery. Monitoring history and silences now persist in named Docker volumes; avoid \texttt{docker compose down -\allowbreak{}v} when preserving evidence.

\begin{quote}\small
\begin{verbatim}
.venv/bin/python scripts/verify_monitoring.py --drill
\end{verbatim}
\end{quote}
This verifies probes, rules, dashboards, Watchdog delivery and a temporary real firing/resolved drill. The TA orchestration runs the non-drill monitoring verification automatically. External notifications are not configured.

\subsection*{Failure investigation}
Check \texttt{reports/\allowbreak{}runtime/\allowbreak{}pipeline.\allowbreak{}log}, stage JSON files, and the matching failed MLflow run. The pipeline console shows the bounded log tail. A diagnostic run named \texttt{diagnostic\_\allowbreak{}failure} deliberately exits with code 3 to verify failure tracking; it is not a failed final training stage. Model-service errors appear in \texttt{docker compose logs model-\allowbreak{}serving backend}.

A missing \texttt{champion} alias requires finishing the pipeline or restoring registration. An offline model cache error requires running the start script with a working public-model download connection. A malformed optional report becomes null in the console and is logged. CPU inference is serialized; do not benchmark concurrency by spawning extra model processes on limited RAM.

\subsection*{Rollback and shutdown}
Review the previous model version's acceptance tags and immutable adapter hash in MLflow. Reset \texttt{champion} to that validated version using \texttt{MlflowClient.\allowbreak{}set\_\allowbreak{}registered\_\allowbreak{}model\_\allowbreak{}alias}, with the local tracking URI set to this repository's SQLite database. Restart \texttt{model-\allowbreak{}serving}, and run \texttt{scripts/\allowbreak{}verify\_\allowbreak{}runtime.\allowbreak{}py} again. Changing the alias alone does not reload an already running model process.

Use \texttt{docker compose stop} to pause the stack or \texttt{docker compose down} to remove its containers/network. The bind-mounted artifacts and MLflow history remain in the project folder. To rerun training on a low-memory machine, stop the model-serving service first and restart it after training completes.

\clearpage
\section*{TA setup and evaluation}
\subsection*{Public resources}
The fresh-checkout script prints these links before doing any work:

\begingroup\small
\begin{longtable}{@{}>{\raggedright\arraybackslash}p{0.3333\dimexpr\textwidth-24pt\relax}>{\raggedright\arraybackslash}p{0.3333\dimexpr\textwidth-24pt\relax}>{\raggedright\arraybackslash}p{0.3333\dimexpr\textwidth-24pt\relax}@{}}
\toprule
\textbf{Resource} & \textbf{Repository} & \textbf{Pinned revision} \\
\midrule
\endfirsthead
\textbf{Resource} & \textbf{Repository} & \textbf{Pinned revision} \\
\midrule
\endhead
Base model and tokenizer & \href{https://huggingface.co/Qwen/Qwen2.5-0.5B-Instruct}{Qwen/Qwen2.5-0.5B-Instruct} & \href{https://huggingface.co/Qwen/Qwen2.5-0.5B-Instruct/tree/7ae557604adf67be50417f59c2c2f167def9a775}{7ae55760\allowbreak{}4adf67be\allowbreak{}50417f59\allowbreak{}c2c2f167\allowbreak{}def9a775} \\
Training/evaluation data & \href{https://huggingface.co/datasets/AI-MO/NuminaMath-1.5}{AI-MO/NuminaMath-1.5} & \href{https://huggingface.co/datasets/AI-MO/NuminaMath-1.5/tree/1b05109f9e5c1ad06c0663519502416c30b300f8}{1b05109f\allowbreak{}9e5c1ad0\allowbreak{}6c066351\allowbreak{}9502416c\allowbreak{}30b300f8} \\
\bottomrule
\end{longtable}
\endgroup
The scripts download the public Qwen weights and tokenizer into the Hugging Face cache. The DVC ingestion stage range-reads an 800-example dataset snapshot; it does not download the entire NuminaMath corpus. The LoRA adapter is built locally. Historical Gemma evidence remains available in the submission bundle, but Gemma is not a dependency of this rebuild. No Hugging Face account or private model repository is required.

\subsection*{Machine prerequisites}
Use Linux with Bash, Python \textbf{3.14.6} including \texttt{venv}, Java \textbf{17 or newer}, Git, \texttt{flock} (util-linux), Docker Engine and the \textbf{Compose v2} plugin. Start the Docker daemon and give the current user permission to run Docker. The scripts check these prerequisites and report missing tools; they do not install system packages or change user permissions.

Allow approximately 8 GiB of available memory and 15 GiB of free disk for the environment, model, Docker images and run artifacts. Internet access is required for first-time packages, images and public resources. Training takes about 25--30 minutes on the demonstration CPU; a fresh dependency/image build adds time and slower machines take longer. Stop other memory-heavy tasks before training.

If needed, choose the interpreter/cache explicitly:

\begin{quote}\small
\begin{verbatim}
export PYTHON_BIN=/absolute/path/to/python3.14
export JAVA_HOME=/absolute/path/to/jdk-17
export MATH_ASSISTANT_HF_CACHE=/absolute/path/to/huggingface
\end{verbatim}
\end{quote}
A project-local \texttt{.\allowbreak{}tools/\allowbreak{}java17} is automatically detected when present. Model downloads and Docker serving use the same cache path. Run the scripts from a terminal; they locate the project root themselves.

\subsection*{Script 1: evaluate the entire workflow}
\begin{quote}\small
\begin{verbatim}
./scripts/ta_evaluate.sh
\end{verbatim}
\end{quote}
Use this after setup when \texttt{.\allowbreak{}venv} already contains the pinned dependencies and the base model is cached. The script performs the following sequence and stops on the first failure:

\begin{enumerate}
\item Check the Python environment and Compose configuration. Acquire a checkout-local evaluation lock and create an individual log.
\item Stop the project's nine services to free training memory, preserving artifacts and volumes.
\item Force the complete DVC graph: environment inspection, dataset ingestion, EDA, Spark cleaning, splitting, CPU benchmark, LoRA training, base evaluation, adapter evaluation, generated-answer acceptance and registry promotion.
\item Ensure a local executable champion exists, restoring registration from verified artifacts only if needed. Update the console's links to the actual local MLflow experiment ID.
\item Run unit tests, display the DVC graph/status, validate submitted acceptance evidence, check Compose/Grafana configuration, regenerate the five-question smoke suite into a separate runtime file, and exercise the MLprojects test entry point.
\item Build the backend/frontend Docker images and recreate all nine services. Recreating model serving loads the current champion.
\item Verify real web/API-to-model inference, MLflow, exporter metrics, Prometheus targets, the provisioned Grafana dashboard, data separation and the completed pipeline. Check that the separate pipeline console responds, then verify the three dashboards, all fifteen scrape targets, eight HTTP probes and recent real Watchdog delivery.
\item Print service health, local links and the pass-result file. Leave the application running for inspection or recording.
\end{enumerate}
Options can be combined:

\begin{quote}\small
\begin{verbatim}
# Display the sequence without downloads, changes, training or service
operations
./scripts/ta_evaluate.sh --dry-run

# Allow DVC to reuse unchanged outputs; rerun missing or changed stages
./scripts/ta_evaluate.sh --cached

# Use existing Docker images for a faster review on the demonstration
machine
./scripts/ta_evaluate.sh --cached --skip-build
\end{verbatim}
\end{quote}
\texttt{-\allowbreak{}-\allowbreak{}cached} can still train if outputs or dependencies changed. \texttt{-\allowbreak{}-\allowbreak{}skip-\allowbreak{}build} requires current images; it is not a fresh-machine setup. The default always retrains and builds. Acceptance checks report a five-case functional smoke score, not broad mathematical accuracy.

\subsection*{Script 2: build a fresh checkout with one command}
\begin{quote}\small
\begin{verbatim}
./scripts/bootstrap_project.sh
\end{verbatim}
\end{quote}
This script prints the pinned model and dataset links, checks machine prerequisites, creates \texttt{.\allowbreak{}venv}, installs \texttt{requirements.\allowbreak{}lock}, downloads the pinned base model/tokenizer, then invokes Script 1. The ingestion stage downloads the bounded dataset, and the rest of the tracked workflow builds the adapter, registration, images and running application. A source-only checkout needs neither the demonstration machine's MLflow database nor its ignored DVC artifact store. An unpacked ZIP without \texttt{.\allowbreak{}git} gets a local Git repository for DVC; the script makes no commits.

For a preview or a repeat setup:

\begin{quote}\small
\begin{verbatim}
./scripts/bootstrap_project.sh --dry-run
./scripts/bootstrap_project.sh --cached
\end{verbatim}
\end{quote}
The selected pipeline rebuilds Qwen outputs; it does not fetch missing historical Gemma artifacts. \texttt{dvc status} may list those standalone historical artifacts as missing in a source-only clone. They are outside the selected rebuild graph and are not used by the deployed model.

\subsection*{Results and local screens}
Each evaluation creates \texttt{reports/\allowbreak{}runtime/\allowbreak{}ta-\allowbreak{}evaluation-\allowbreak{}<UTC timestamp>-\allowbreak{}<PID>.\allowbreak{}log} and a matching JSON result with \texttt{passed} or \texttt{failed}, the last phase and completed steps. Bootstrap adds a \texttt{bootstrap-\allowbreak{}<UTC timestamp>-\allowbreak{}<PID>.\allowbreak{}log}. Failures return a nonzero exit code; a failed pipeline prevents startup and a failed runtime check prevents a pass result. Preflight errors print directly before a run log is created.

Detailed pipeline output is in \texttt{reports/\allowbreak{}runtime/\allowbreak{}pipeline.\allowbreak{}log}. MLflow stores parent/stage runs and artifacts in the checkout's SQLite-backed tracking store. Unit results are in \texttt{reports/\allowbreak{}test-\allowbreak{}results.\allowbreak{}xml}; fresh smoke results are in \texttt{reports/\allowbreak{}runtime/\allowbreak{}qwen\_\allowbreak{}generation\_\allowbreak{}rerun.\allowbreak{}json}; live checks are in \texttt{reports/\allowbreak{}runtime/\allowbreak{}runtime-\allowbreak{}verification.\allowbreak{}json}. Cached stages retain their earlier evidence rather than acquiring invented new measurements.

\begingroup\small
\begin{longtable}{@{}>{\raggedright\arraybackslash}p{0.5000\dimexpr\textwidth-12pt\relax}>{\raggedright\arraybackslash}p{0.5000\dimexpr\textwidth-12pt\relax}@{}}
\toprule
\textbf{Screen} & \textbf{Local URL} \\
\midrule
\endfirsthead
\textbf{Screen} & \textbf{Local URL} \\
\midrule
\endhead
Math workspace & \url{http://localhost:8080/} \\
Pipeline console & \url{http://localhost:8081/} \\
MLflow experiments and registry & \url{http://localhost:5000/} \\
FastAPI documentation & \url{http://localhost:8000/docs} \\
MLflow model health & \url{http://localhost:5002/ping} \\
Prometheus & \url{http://localhost:19090/} \\
Grafana operations & \url{http://localhost:13000/d/local-math-assistant} \\
Grafana Alert Center & \url{http://localhost:13000/d/math-alert-center} \\
Grafana ML lifecycle & \url{http://localhost:13000/d/math-mlops} \\
Alertmanager and silences & \url{http://localhost:19093/} \\
Local alert inbox & \url{http://localhost:19199/} \\
\bottomrule
\end{longtable}
\endgroup
On failure, inspect the reported phase and log, then use \texttt{docker compose logs model-\allowbreak{}serving backend} if startup/inference failed. Services remain in their last state for diagnosis; a failure after the stop step may leave them stopped. Fix the reported issue and rerun, using \texttt{-\allowbreak{}-\allowbreak{}cached} when appropriate. Stop the stack with \texttt{docker compose stop} after evaluation.

Both orchestration scripts work locally, preserve prior MLflow runs, and perform no Git commit, GitHub push, cloud deployment or artifact upload.

\clearpage
\section*{Evaluation Rubric Evidence Map}
This maps the supplied evaluation rubric to concrete repository files and demonstration evidence. The technical categories total 20 points; the viva contributes a further 5 points through the presenter's own explanations.

\begingroup\small
\begin{longtable}{@{}>{\raggedright\arraybackslash}p{0.3333\dimexpr\textwidth-24pt\relax}>{\raggedright\arraybackslash}p{0.3333\dimexpr\textwidth-24pt\relax}>{\raggedright\arraybackslash}p{0.3333\dimexpr\textwidth-24pt\relax}@{}}
\toprule
\textbf{Rubric category} & \textbf{Points} & \textbf{Implementation and evidence} \\
\midrule
\endfirsthead
\textbf{Rubric category} & \textbf{Points} & \textbf{Implementation and evidence} \\
\midrule
\endhead
Web frontend UI/UX & 3 & Independent responsive workspace, example questions, keyboard shortcut, input validation, progress/error feedback, safe answer rendering, copy action and in-app guide. \texttt{frontend/\allowbreak{}}, \texttt{docs/\allowbreak{}user\_\allowbreak{}manual.\allowbreak{}md}, screenshots and browser test results. \\
ML pipeline visualization & 2 & Separate \texttt{pipeline.\allowbreak{}html} console with eleven stages, data counts, durations, statuses, run links, log tail and operations tool links. \texttt{GET /\allowbreak{}v1/\allowbreak{}pipeline/\allowbreak{}status}, \texttt{/\allowbreak{}v1/\allowbreak{}pipeline/\allowbreak{}logs}, MLflow parent/child runs and DVC DAG. \\
Design principles & 2 & Architecture SVG/Mermaid and block explanation; HLD rationale/tradeoffs; LLD endpoint inputs/outputs/errors. Frontend  ->  configurable REST  ->  FastAPI  ->  separate MLflow inference. \\
Implementation & 2 & Typed models/configuration, narrow model-service object, functions for pipelines, exception boundaries, bounded logging labels, source documentation and automated tests. \\
Testing & 1 & \texttt{docs/\allowbreak{}test\_\allowbreak{}plan.\allowbreak{}md}, test case inventory and acceptance criteria; \texttt{docs/\allowbreak{}test\_\allowbreak{}report.\allowbreak{}md}; \texttt{reports/\allowbreak{}test-\allowbreak{}results.\allowbreak{}xml}; live runtime verification and browser checks. \\
Data engineering & 2 & Pinned bounded Parquet ingestion, local PySpark cleaning/normalization/deduplication, deterministic split with counts/hashes, EDA and measured throughput. \texttt{src/\allowbreak{}data/\allowbreak{}}, DVC stages and quality report. \\
Source control and CI & 2 & Existing Git repository; DVC DAG, dependency tracking, lock file and local artifact remote. Reproducible runner and evaluation script perform the local integration/gating sequence. Source/model/data/config identities are preserved. No remote CI run or GitHub upload is claimed. \\
Experiment tracking & 2 & MLflow parent/stage runs; trainer step loss; timings, source snapshot, environment freeze, Git dirty identity, output hashes and explicit metrics/artifacts beyond autologging; verified model versions and alias. \\
Exporter instrumentation and visualization & 2 & API, data pipeline, model exporter and process metrics; Prometheus scrapes; provisioned Grafana datasource/dashboard; seventeen alert rules, six recording rules, eight HTTP probes, Alertmanager, a persistent local inbox and a verified firing/resolved drill. Runtime checks record target health and generated metrics. \\
Software packaging & 2 & Real MLflow Models Serve \texttt{/\allowbreak{}invocations}, MLproject/python environment, FastAPI contract, independent backend/frontend images, nine-service Docker Compose and successful live HTTP boundary checks. \\
Viva & 5 & \texttt{docs/\allowbreak{}viva\_\allowbreak{}guide.\allowbreak{}md} supports the project/tools/design discussion. The presenter must demonstrate their own understanding and review Assignments 1-10. Those personal viva outcomes cannot be verified by repository automation. \\
\bottomrule
\end{longtable}
\endgroup
\subsection*{Mandatory document locations}
\begin{enumerate}
\item Architecture diagram and block explanation: \texttt{docs/\allowbreak{}architecture.\allowbreak{}svg}, \texttt{docs/\allowbreak{}architecture.\allowbreak{}md}.
\item High-level design with rationale: \texttt{docs/\allowbreak{}high\_\allowbreak{}level\_\allowbreak{}design.\allowbreak{}md}.
\item Low-level design and API I/O: \texttt{docs/\allowbreak{}low\_\allowbreak{}level\_\allowbreak{}design.\allowbreak{}md}, live OpenAPI.
\item Test plan/cases: \texttt{docs/\allowbreak{}test\_\allowbreak{}plan.\allowbreak{}md}; measured report: \texttt{docs/\allowbreak{}test\_\allowbreak{}report.\allowbreak{}md}.
\item Nontechnical user manual: \texttt{docs/\allowbreak{}user\_\allowbreak{}manual.\allowbreak{}md}, \texttt{frontend/\allowbreak{}help.\allowbreak{}html}.
\end{enumerate}
\texttt{docs/\allowbreak{}project\_\allowbreak{}report.\allowbreak{}pdf} consolidates the documents for review. \texttt{docs/\allowbreak{}recording\_\allowbreak{}guide.\allowbreak{}md} supplies a suggested video sequence. The local submission ZIP includes code, required documents, reports, artifact objects, a file manifest and archive checksum.

\subsection*{Evidence boundaries}
Historical Gemma comparisons remain identifiable as earlier experiments. The fresh selected-model run uses Qwen. The generation smoke score is limited to five declared questions. Source is intentionally left uncommitted for the owner to review and commit; runs record the existing commit plus dirty state and preserve source snapshots. Artifact restoration is supported by the archive, while a source-only future GitHub clone requires the separate artifact bundle or a full rebuild.

\clearpage
\section*{Demonstration Recording Guide}
Start the stack with \texttt{.\allowbreak{}/\allowbreak{}scripts/\allowbreak{}start\_\allowbreak{}demo.\allowbreak{}sh}. Keep the project terminal and the following local pages ready. Record the workflow and explain each step as you show it.

\begingroup\small
\begin{longtable}{@{}>{\raggedright\arraybackslash}p{0.3333\dimexpr\textwidth-24pt\relax}>{\raggedright\arraybackslash}p{0.3333\dimexpr\textwidth-24pt\relax}>{\raggedright\arraybackslash}p{0.3333\dimexpr\textwidth-24pt\relax}@{}}
\toprule
\textbf{Segment} & \textbf{Show} & \textbf{Explain} \\
\midrule
\endfirsthead
\textbf{Segment} & \textbf{Show} & \textbf{Explain} \\
\midrule
\endhead
0:00-0:40 & Workspace at localhost:8080 & The user problem, local CPU inference, question examples and navigation \\
0:40-1:30 & Submit an empty question, then solve \texttt{2x + 3 = 11} & Clear input feedback, generated result, response time, copy function \\
1:30-2:20 & Fraction example; resize the window & Intuitive workflow and responsive layout \\
2:20-3:40 & Pipeline screen & Cleaned/split counts, actual timings, successful stages and execution log \\
3:40-4:40 & \texttt{dvc dag}, \texttt{dvc status}, \texttt{dvc.\allowbreak{}yaml}, \texttt{dvc.\allowbreak{}lock} & Dependencies, versioned data/model artifacts and reproducible commands \\
4:40-5:50 & MLflow mathematics-pipeline experiment & Parent/child runs, training loss curve, parameters, source snapshot and output artifacts \\
5:50-6:40 & MLflow registry & Verified candidate version, acceptance tag, hash and \texttt{champion} alias; show a historical failed diagnostic run separately \\
6:40-7:40 & Grafana dashboard and Prometheus targets & Live request rate/latency, pipeline durations, Spark throughput, model metric and all scrape targets up \\
7:40-8:30 & FastAPI \texttt{/\allowbreak{}docs}, Compose services & Loose coupling: Nginx  ->  FastAPI  ->  MLflow model serving; distinct containers \\
8:30-9:20 & Test report and mandatory design/manual documents & Test coverage, acceptance scope, model limitations and design rationale \\
\bottomrule
\end{longtable}
\endgroup
For an additional progress clip, run \texttt{.\allowbreak{}venv/\allowbreak{}bin/\allowbreak{}python -\allowbreak{}m src.\allowbreak{}pipeline.\allowbreak{}run -\allowbreak{}-\allowbreak{}target split -\allowbreak{}-\allowbreak{}force} and show the console/MLflow update. This is the real data pipeline and takes seconds, including Spark startup. For training footage, use a full \texttt{-\allowbreak{}-\allowbreak{}force} run with the inference service stopped; capture the live training curve and terminal progress, then restart the stack afterward. Do not describe a cached DVC run as new training.

The five-case generation suite is a bounded smoke test. Explain its scope clearly when showing 100\%: it means five accepted outputs, not perfect mathematics performance. Distinguish held-out completion loss from correctness.

\subsection*{Alerting demonstration}
Show the service-operations overview, Alert Center and ML lifecycle dashboard. Open the local inbox at \url{http://localhost:19199/} to show Watchdog delivery. Run \texttt{.\allowbreak{}venv/\allowbreak{}bin/\allowbreak{}python scripts/\allowbreak{}verify\_\allowbreak{}monitoring.\allowbreak{}py -\allowbreak{}-\allowbreak{}drill} and record the clearly labelled firing and resolved drill notifications. It typically takes one to two minutes. Explain grouping, hold periods, traffic gates and timed silences; the drill is a synthetic delivery test, not a real outage. No application service is stopped.

\clearpage
\section*{Viva Preparation}
\subsection*{Explain the project in one minute}
MathStudio is a local mathematics practice assistant with an independently deployed web UI and REST backend. A fine-tuned Qwen LoRA adapter is exposed through MLflow Models Serve. DVC orchestrates ingestion, Spark cleaning, deterministic splitting, resource benchmarking, training, evaluation and registration. MLflow records experiments and a verified champion decision. Prometheus and Grafana show API, model, data-pipeline and process behavior. The pipeline screen combines lifecycle evidence and tool links.

\subsection*{Questions to practise}
\begingroup\small
\begin{longtable}{@{}>{\raggedright\arraybackslash}p{0.5000\dimexpr\textwidth-12pt\relax}>{\raggedright\arraybackslash}p{0.5000\dimexpr\textwidth-12pt\relax}@{}}
\toprule
\textbf{Question} & \textbf{Points to explain} \\
\midrule
\endfirsthead
\textbf{Question} & \textbf{Points to explain} \\
\midrule
\endhead
Why use a small model and LoRA? & CPU/RAM constraints, limited trainable parameters, adapter portability, capability tradeoffs \\
Why Spark for only 800 rows? & Demonstrates the required data-engineering engine; bounded CPU experiment; overhead and scaling tradeoff \\
What do DVC and Git each track? & Git tracks code/configuration and object pointers; DVC tracks artifact contents and dependency identities \\
What does MLflow add? & Nested experiment runs, live metrics, explicit parameters/artifacts, executable pyfunc models and registry aliases \\
How is MLflow actually used for inference? & FastAPI calls \texttt{/\allowbreak{}invocations} on a distinct MLflow Models Serve process using a registered model signature \\
How do you prevent leakage? & Split normalized questions with a deterministic hash, keep question duplicates in one partition, score held-out test rows \\
Why completion-only labels? & Ignore system/user prompt tokens; optimize the response; count fully truncated completions \\
Does lower loss prove correct mathematics? & No; loss measures fit to reference completions; generated-answer tests have a separate declared threshold \\
Is the 100\% score general accuracy? & No; it is limited to the five-case smoke suite and must be presented with that sample size \\
How is promotion controlled? & Recompute case decisions, verify threshold/base revision/adapter SHA-256, retain failed candidates, update champion only on acceptance \\
How is inference protected against repeated requests? & Serialized model lock, UI duplicate-submit guard, bounded input/output and timeouts \\
How are failures observed? & Stage failure JSON and failed MLflow runs; stable HTTP errors; metrics and local alert rules \\
How do you reproduce an uncommitted run? & Git commit plus dirty flag, saved source snapshot, pip freeze, pinned revisions, output hashes and DVC lock \\
How would you roll back? & Verify a prior accepted registry version, reset alias, restart model serving and repeat live checks \\
What needs improvement before production? & Broader unseen evaluation, real drift baselines, authentication/TLS, multi-user resource management and scalable experiment storage \\
\bottomrule
\end{longtable}
\endgroup
Review your own Assignments 1-10 for the viva criterion specific to those assignments. Their complete question papers and your answers are not part of this repository, so this guide does not claim to cover them. Be able to open each relevant source file and explain the implementation without relying on memorized wording.

\end{document}
